% ===========================================================================
% Poste 12 — Crédit remboursable pour frais médicaux (Québec)
% ===========================================================================
\poste{12}{Crédit d'impôt remboursable pour frais médicaux (Québec)}{QC\_medic}

\pdesc Crédit remboursable destiné aux travailleurs à faible revenu dont les
frais médicaux dépassent un seuil relatif à leur revenu (ligne 462 de la
déclaration). \textbf{Dans ce modèle, ce poste vaut toujours zéro} --- il
est documenté ici pour expliquer pourquoi.

\pobj Alléger les frais médicaux des travailleurs à faible revenu, même
lorsqu'ils ne paient pas assez d'impôt pour profiter du crédit non
remboursable usuel.

\pregle La structure officielle du crédit, telle que codée :
\begin{enumerate}
  \item admissibilité : un revenu de travail d'au moins $r_{min}$ ;
  \item frais admissibles : la part des frais médicaux $F$ qui excède
        \pc{3} du revenu familial net,
        $F_a = \pos{F - 0{,}03 \times \RFN}$ ;
  \item crédit avant réduction : $c = \min(0{,}25 \times F_a,\; c_{max})$ ;
  \item réduction selon le revenu : l'excédent $e = \pos{\RFN - s}$ est
        soustrait, au taux de \pc{5}.
\end{enumerate}

\pnotes Deux particularités, reproduites par fidélité au calculateur
officiel, expliquent le résultat systématiquement nul.
\begin{itemize}
  \item \emph{La seule dépense médicale du modèle est la prime RAMQ}
        (poste~5) : aucune autre entrée de frais médicaux n'existe. Or la
        prime, plafonnée à environ \D{750}, ne dépasse \emph{jamais} \pc{3}
        du revenu familial net au-dessus duquel elle devient exigible : les
        frais admissibles $F_a$ sont donc toujours nuls, et le crédit avec
        eux. Ce résultat a été vérifié empiriquement sur l'ensemble de la
        grille de validation ;
  \item \emph{la forme de la réduction du calculateur officiel est
        anormale} : il calcule $\pos{(c - e) \times 0{,}05}$ là où la règle
        réelle est $\pos{c - 0{,}05 \times e}$. Comme $c$ est toujours nul,
        l'anomalie ne se manifeste jamais ; le code la reproduit telle
        quelle, avec ce commentaire, plutôt que de \og corriger \fg{}
        silencieusement la référence.
\end{itemize}

\pparams

\begin{center}
\small
\begin{tabular}{l r r}
\toprule
Paramètre & 2025 & 2026 \\
\midrule
Taux du crédit                       & \pc{25}   & \pc{25} \\
Crédit maximal ($c_{max}$)           & \D{1466}  & \D{1496} \\
Revenu de travail minimal ($r_{min}$)& \D{3750}  & \D{3825} \\
Seuil de réduction ($s$)             & \D{28335} & \D{28915} \\
Taux de réduction                    & \pc{5}    & \pc{5} \\
Plancher des frais (part du revenu)  & \pc{3}    & \pc{3} \\
\bottomrule
\end{tabular}
\end{center}

\prefs Loi sur les impôts (RLRQ, c.~I-3) ; déclaration TP-1, ligne 462 ;
Revenu Québec ; ministère des Finances du Québec, \emph{Paramètres du
régime d'imposition}.
